% !TEX TS-program = pdflatex
% !TEX encoding = UTF-8 Unicode

% This file is a template using the "beamer" package to create slides for a talk or presentation
% - Giving a talk on some subject.
% - The talk is between 15min and 45min long.
% - Style is ornate.

% MODIFIED by Jonathan Kew, 2008-07-06
% The header comments and encoding in this file were modified for inclusion with TeXworks.
% The content is otherwise unchanged from the original distributed with the beamer package.

\documentclass{beamer}
\setbeamercovered{invisible}



% Copyright 2004 by Till Tantau <tantau@users.sourceforge.net>.
%
% In principle, this file can be redistributed and/or modified under
% the terms of the GNU Public License, version 2.
%
% However, this file is supposed to be a template to be modified
% for your own needs. For this reason, if you use this file as a
% template and not specifically distribute it as part of a another
% package/program, I grant the extra permission to freely copy and
% modify this file as you see fit and even to delete this copyright
% notice. 


\mode<presentation>
{
  \usetheme{Warsaw}
  % or ...

  % or whatever (possibly just delete it)
}


\usepackage[english]{babel}
% or whatever

\usepackage[utf8]{inputenc}
\usepackage{qrcode}
\usepackage{hyperref}
\usepackage{ytableau}
\usepackage{longtable}
% or whatever

\usepackage{times}
\usepackage[T1]{fontenc}
% Or whatever. Note that the encoding and the font should match. If T1
% does not look nice, try deleting the line with the fontenc.

%%% MATH RELATED

\title{MATH 180 Strategies of Problem Solving} 
\subtitle
{} % (optional)

\author[W.R. Casper] % (optional, use only with lots of authors)
{W.R. Casper}
% - Use the \inst{?} command only if the authors have different
%   affiliation.

\institute[California State University Fullerton] % (optional, but mostly needed)
{
  Department of Mathematics\\
  California State University Fullerton}
% - Use the \inst command only if there are several affiliations.
% - Keep it simple, no one is interested in your street address.

\subject{Talks}
% This is only inserted into the PDF information catalog. Can be left
% out. 



% If you have a file called "university-logo-filename.xxx", where xxx
% is a graphic format that can be processed by latex or pdflatex,
% resp., then you can add a logo as follows:

% \pgfdeclareimage[height=0.5cm]{university-logo}{university-logo-filename}
% \logo{\pgfuseimage{university-logo}}



% Delete this, if you do not want the table of contents to pop up at
% the beginning of each subsection:
\AtBeginSubsection[]
{
  \begin{frame}<beamer>{Outline}
    \tableofcontents[currentsection,currentsubsection]
  \end{frame}
}


% If you wish to uncover everything in a step-wise fashion, uncomment
% the following command: 

%\beamerdefaultoverlayspecification{<+->}


\begin{document}

\begin{frame}
  \titlepage
\end{frame}

\begin{frame}{Outline}
  \tableofcontents
  % You might wish to add the option [pausesections]
\end{frame}

% Since this a solution template for a generic talk, very little can
% be said about how it should be structured. However, the talk length
% of between 15min and 45min and the theme suggest that you stick to
% the following rules:  

% - Exactly two or three sections (other than the summary).
% - At *most* three subsections per section.
% - Talk about 30s to 2min per frame. So there should be between about
%   15 and 30 frames, all told.

\section{SoPS Lecture 8}
\subsection{Break the Cube}


\begin{frame}{Review}
  \begin{block}{Last time}
    \begin{itemize}
\pause
      \item we explored Battleship firing algorithms
\pause
      \item we talked about static vs. dynamic algorithms
\pause
      \item we found good patterns for static algorithms
    \end{itemize}
  \end{block}
\end{frame}

\begin{frame}{Break the Cube}
  \begin{block}{Activity (60 mins)}
    \begin{itemize}
\pause
      \item split into groups of up to $4$
\pause
      \item take some time to read the rules
\pause
      \item play a few games
\pause
      \item {\color{red}record data:} how many clues until a correct guess
    \end{itemize}
  \end{block}
\end{frame}
    
\begin{frame}{Break the Cube: Debrief}
\pause
\textbf{Let's combine the data ...}
\pause
  \begin{block}{Question}
    How many clues do we need to find a solution?
  \end{block}
  \begin{block}{Question}
    What kinds of boards took the most clues?
  \end{block}
%\pause
%  \begin{block}{Submit your conjecture:}
%    \begin{center}
%      \qrcode[hyperlink,height=2cm]{https://padlet.com/williamrileycasper/conjectures-for-losing-positions-33ai54j4kifc9hx6}
%    \end{center}
%  \end{block}
%\pause
%  \begin{center}
%    There seems to be some mysterious math behind the scenes...
%  \end{center}
\end{frame}

\begin{frame}{Break the Cube: Solitaire}
  \begin{block}{Problem}
    Suppose we have the clues:
    \begin{itemize}
      \item B: green, blank, blank
      \item C: yellow, gray, blank
      \item J: green, blank, blank
    \end{itemize}
    with gray=small, green=medium, and yellow=large.\\
    Find the block setup.
  \end{block}
\end{frame}

\begin{frame}{Prove it!}
    \pause
  \begin{block}{Proof}
    A proof is a logical argument that explains why a statement must be true no matter what.
  \end{block}
    \pause
  \begin{block}{Example:}
    For any integer $n$, the number $n^2+n$ is even.\\
    \pause
    \textbf{Proof:}\\
    If $n$ is an integer, either $n$ or $n+1$ will be even.\\
    Therefore $n^2+n = n(n+1)$ is the product of an even number with another number.\\
    Any number times an even number is even.\\
    Therefore $n^2+n$ must be even.
  \end{block}
\end{frame}


\begin{frame}{Build your own}
  \begin{block}{Activity}
    Create a Break the Cube solitaire problem, meaning:
    \begin{itemize}
      \item specify the colors of the small, medium, large blocks
      \item enough clues to determine a \textbf{unique} setup
    \end{itemize}
  \end{block}
  \pause
  \begin{block}{Activity}
    Swap your puzzle with another group.
    \begin{itemize}
      \item find a solution of their puzzle
      \item explain why the solution is unique (or find another)
      \item try to write a proof
    \end{itemize}
  \end{block}
\end{frame}

\subsection{Wrapping up}

\begin{frame}{Discussion}
\pause
  \begin{block}{Question}
    What makes a good choice of setup for your Break the Cube board?
  \end{block}
\pause
  \begin{block}{Question}
    What does it take to \emph{prove} that given a set of clues your solution is right?
  \end{block}
\end{frame}


\begin{frame}{Question of the Day}
\pause
  \begin{block}{Puzzle}
    Suppose we wanted to try to count the number of possible arrangements of the three blocks.  How might we begin to go about it?
  \end{block}
\end{frame}

\begin{frame}{Final thoughts}
  \begin{block}{Reflection}
    \begin{itemize}
\pause
      \item Describe what a proof is in your own words.
\pause
      \item Find an example of a proof online and write about it.  What is it proving?  What is the crux of the argument, in your own words?
    \end{itemize}
  \end{block}
\begin{center}
Submit your responses by \textbf{Tonight!}
\end{center}
\end{frame}


\end{document}


